%==============================================================================
\section{Act III: The Hero (\texttt{Tracking.h/.cc})}
%==============================================================================

The \texttt{Tracking} class is where the magic happens. It takes raw images and produces camera poses.

\subsection{The Tracking States}

Tracking operates like a state machine:

\begin{lstlisting}[caption={Tracking states}]
enum eTrackingState {
    SYSTEM_NOT_READY = -1,  // Still initializing
    NO_IMAGES_YET = 0,      // Haven't received any images
    NOT_INITIALIZED = 1,    // Trying to initialize the map
    OK = 2,                 // Normal operation
    RECENTLY_LOST = 3,      // Just lost tracking (trying to recover)
    LOST = 4,               // Completely lost
    OK_KLT = 5              // Tracking with KLT (fallback)
};
\end{lstlisting}

\textbf{The lifecycle:}
\begin{enumerate}
    \item Start in \texttt{NO\_IMAGES\_YET}
    \item First image $\rightarrow$ transition to \texttt{NOT\_INITIALIZED}
    \item After successful initialization $\rightarrow$ transition to \texttt{OK}
    \item If tracking fails $\rightarrow$ \texttt{RECENTLY\_LOST} $\rightarrow$ try relocalization
    \item If relocalization fails $\rightarrow$ \texttt{LOST}
\end{enumerate}

\subsection{The Main Entry: GrabImageMonocular}

\texttt{GrabImageMonocular} is called by System for each new image:

\begin{lstlisting}[caption={GrabImageMonocular simplified}]
Sophus::SE3f Tracking::GrabImageMonocular(const cv::Mat &im, 
                                          const double &timestamp, 
                                          string filename)
{
    // 1. Store the image
    mImGray = im;
    
    // 2. Create a Frame (feature extraction happens here!)
    mCurrentFrame = Frame(mImGray, timestamp, mpORBextractorLeft, 
                          mpORBVocabulary, mpCamera, mDistCoef, 
                          mbf, mThDepth, &mLastFrame, mImuCalib);
    
    // 3. The main tracking logic
    Track();
    
    // 4. Return the estimated pose
    return mCurrentFrame.GetPose();
}
\end{lstlisting}

\subsection{Inside \texttt{Track()}: The Heart of SLAM}

The \texttt{Track()} function is a 500+ line behemoth, but its structure is simple:

\begin{lstlisting}[caption={Track() simplified structure}]
void Tracking::Track()
{
    // PHASE 1: INITIALIZATION
    if(mState == NOT_INITIALIZED)
    {
        MonocularInitialization();  // Try to initialize the map
        return;
    }
    
    // PHASE 2: POSE ESTIMATION (if initialized)
    bool bOK = false;
    
    // First attempt: Use motion model
    if(mbVelocity)
        bOK = TrackWithMotionModel();
    
    // Fallback: Match against reference keyframe
    if(!bOK)
        bOK = TrackReferenceKeyFrame();
    
    // If still failing: Try relocalization
    if(!bOK)
        bOK = Relocalization();
    
    // PHASE 3: TRACK LOCAL MAP (refine pose using more points)
    if(bOK)
        bOK = TrackLocalMap();
    
    // PHASE 4: KEYFRAME DECISION
    if(bOK && NeedNewKeyFrame())
        CreateNewKeyFrame();
    
    // Update state
    mState = bOK ? OK : LOST;
}
\end{lstlisting}

\subsection{Phase 1: Monocular Initialization}

Monocular SLAM has a chicken-and-egg problem:
\begin{itemize}
    \item To estimate the camera pose, we need 3D points
    \item To triangulate 3D points, we need two camera poses
\end{itemize}

The solution: \textbf{Two-View Reconstruction}

\begin{lstlisting}[caption={MonocularInitialization simplified}]
void Tracking::MonocularInitialization()
{
    // Step 1: First frame becomes the "initial frame"
    if(mState == NO_IMAGES_YET)
    {
        mInitialFrame = Frame(mCurrentFrame);
        return;
    }
    
    // Step 2: Match features between initial and current frame
    vector<int> vMatches;
    int nMatches = ORBmatcher::SearchForInitialization(
        mInitialFrame, mCurrentFrame, vMatches);
    
    if(nMatches < 100)
        return;  // Not enough matches, wait for better frame
    
    // Step 3: Compute Homography and Fundamental matrix
    // (two competing geometric models)
    cv::Mat H, F;
    float SH, SF;  // Scores for each model
    ComputeH21(vMatches, SH, H);
    ComputeF21(vMatches, SF, F);
    
    // Step 4: Choose the better model and recover pose
    cv::Mat R, t;
    vector<cv::Point3f> vP3D;  // Triangulated 3D points
    
    float RH = SH / (SH + SF);  // "Homographyness"
    
    if(RH > 0.4)
        ReconstructH(H, R, t, vP3D);  // Planar scene
    else
        ReconstructF(F, R, t, vP3D);  // General scene
    
    // Step 5: Create the initial map
    CreateInitialMapMonocular();
    
    mState = OK;  // Initialization complete!
}
\end{lstlisting}

\textbf{Key insight:} After initialization, we have:
\begin{itemize}
    \item Two KeyFrames (the initial frame and current frame)
    \item A set of 3D MapPoints (triangulated from the two views)
    \item An arbitrary scale (monocular SLAM cannot recover absolute scale)
\end{itemize}

\subsection{Phase 2: Pose Estimation}

Once initialized, we estimate the camera pose for each new frame.

\subsubsection{TrackWithMotionModel}

This is the fastest method. We assume the camera is moving at constant velocity:

\begin{lstlisting}[caption={TrackWithMotionModel concept}]
bool Tracking::TrackWithMotionModel()
{
    // 1. Predict pose using velocity
    Sophus::SE3f Tcw_predicted = mVelocity * mLastFrame.GetPose();
    mCurrentFrame.SetPose(Tcw_predicted);
    
    // 2. Project MapPoints from last frame into current frame
    // 3. Search for matches near projected locations
    int nMatches = ORBmatcher::SearchByProjection(mCurrentFrame, mLastFrame);
    
    if(nMatches < 20)
        return false;
    
    // 4. Optimize pose using matched points
    Optimizer::PoseOptimization(&mCurrentFrame);
    
    return true;
}
\end{lstlisting}

\subsubsection{TrackReferenceKeyFrame}

Fallback when motion model fails:

\begin{lstlisting}[caption={TrackReferenceKeyFrame concept}]
bool Tracking::TrackReferenceKeyFrame()
{
    // 1. Compute Bag-of-Words for current frame
    mCurrentFrame.ComputeBoW();
    
    // 2. Match against reference keyframe using BoW
    int nMatches = ORBmatcher::SearchByBoW(mpReferenceKF, mCurrentFrame);
    
    if(nMatches < 15)
        return false;
    
    // 3. Use reference keyframe pose as initial guess
    mCurrentFrame.SetPose(mpReferenceKF->GetPose());
    
    // 4. Optimize pose
    Optimizer::PoseOptimization(&mCurrentFrame);
    
    return true;
}
\end{lstlisting}

\subsection{Phase 3: Track Local Map}

After initial pose estimation, we refine it using more MapPoints:

\begin{lstlisting}[caption={TrackLocalMap concept}]
bool Tracking::TrackLocalMap()
{
    // 1. Find all keyframes that share MapPoints with current frame
    UpdateLocalKeyFrames();
    
    // 2. Collect all MapPoints visible from those keyframes
    UpdateLocalPoints();  // mvpLocalMapPoints
    
    // 3. Project local MapPoints into current frame
    //    Search for matches
    SearchLocalPoints();
    
    // 4. Final pose optimization with all matches
    Optimizer::PoseOptimization(&mCurrentFrame);
    
    // 5. Count inliers
    mnMatchesInliers = CountInliers();
    
    return (mnMatchesInliers > 30);
}
\end{lstlisting}

\textbf{Key insight:} Local Map Tracking uses many more points than the initial pose estimation, giving a more accurate and robust pose.

\subsection{Phase 4: KeyFrame Decision}

Not every frame becomes a KeyFrame. We only insert new KeyFrames when:

\begin{lstlisting}[caption={NeedNewKeyFrame simplified}]
bool Tracking::NeedNewKeyFrame()
{
    // 1. Minimum time since last keyframe
    if(mCurrentFrame.mnId < mnLastKeyFrameId + mMinFrames)
        return false;
    
    // 2. Local mapping must not be busy
    if(mpLocalMapper->isStopped() || mpLocalMapper->stopRequested())
        return false;
    
    // 3. Significant change in observation
    float refRatio = 0.9f;  // 90% of reference KF points still visible
    int nTrackedClose = CountClosePoints();
    
    bool c1 = (mnMatchesInliers < mpReferenceKF->TrackedMapPoints * refRatio);
    bool c2 = (nTrackedClose < 100);  // Few close points
    
    return (c1 || c2);
}
\end{lstlisting}

\textbf{Key insight:} KeyFrames are inserted sparingly to avoid redundancy, but frequently enough to maintain map quality.

\subsection{The Key Member Variables}

\begin{lstlisting}[caption={Important Tracking members}]
// The two most important frames
Frame mCurrentFrame;   // The frame being processed NOW
Frame mLastFrame;      // The previous frame

// Tracking state
eTrackingState mState;

// Current matches
int mnMatchesInliers;

// Local map
KeyFrame* mpReferenceKF;
vector<KeyFrame*> mvpLocalKeyFrames;
vector<MapPoint*> mvpLocalMapPoints;

// Motion model
Sophus::SE3f mVelocity;
bool mbVelocity;

// Feature extraction
ORBextractor* mpORBextractorLeft;
ONNXFeatureExtractor* mpXFeatExtractor;  // Your XFeat mode
\end{lstlisting}
